\documentclass[11pt,a4paper]{article}
\usepackage[T1]{fontenc}
\usepackage{lmodern,microtype}
\usepackage[margin=27mm,headheight=15pt]{geometry}
\usepackage{amsmath,amssymb,amsthm,mathtools}
\usepackage{xcolor,booktabs,tabularx,longtable,enumitem}
\definecolor{inkblue}{HTML}{244860}
\usepackage[colorlinks=true,linkcolor=inkblue,urlcolor=inkblue]{hyperref}
\usepackage{fancyhdr}
\pagestyle{fancy}\fancyhf{}
\fancyhead[L]{\small Exact symplectic comparison}
\fancyhead[R]{\small Weak taming}
\fancyfoot[C]{\small\thepage}
\setlength{\parskip}{4pt}
\setlength{\parindent}{0pt}
\setlength{\emergencystretch}{2em}
\setlist{itemsep=3pt,topsep=5pt}
\newtheorem{theorem}{Theorem}[section]
\newtheorem{lemma}[theorem]{Lemma}
\newtheorem{proposition}[theorem]{Proposition}
\newtheorem{corollary}[theorem]{Corollary}
\theoremstyle{definition}
\newtheorem{definition}[theorem]{Definition}
\newtheorem{remark}[theorem]{Remark}
\newcommand{\R}{\mathbb R}
\newcommand{\C}{\mathbb C}
\newcommand{\dd}{\mathrm d}
\newcommand{\id}{\mathrm{id}}
\newcommand{\norm}[1]{\left\lVert#1\right\rVert}
\newcommand{\abs}[1]{\left|#1\right|}
\newcommand{\ip}[2]{\left\langle#1,#2\right\rangle}
\newcommand{\flatmap}{\omega^\flat}
\newcommand{\dual}{g^*}
\newcommand{\sref}{\mathcal R}
\allowdisplaybreaks[2]

\hypersetup{pdftitle={Exact symplectic comparison: mathematical proof blueprint}}
\title{Exact symplectic comparison under weak taming\\[5pt]
\large A mathematical proof blueprint}
\author{Exact Liouville comparison project}
\date{24 October 2026}
\begin{document}
\maketitle
\begin{abstract}
We give the four-path Moser proof of the radial exact comparison
theorem. The argument constructs a small radial coefficient from the
inverse-form growth bound and a large coefficient dominating an
otherwise unrestricted primitive. A scalar exhaustion controls the
first and third flows; two different barriers control the forward
and reverse directions of the interpolation flows. We then prove the
exact pullback identity with its smooth integral gauge and record
the composition order explicitly. We also construct a smooth
cylindrical compatible structure and extend arbitrary prescribed
compatible contact data, preserving the entire closed end. The final sections derive the
finite-type Liouville, properly embedded Stein, and common-reference
applications. All estimates use the stated reference metric.
\end{abstract}

\section{The theorem and the architecture of its proof}
Let $X$ be a Hausdorff second-countable smooth real manifold without
boundary, of one fixed arbitrary dimension $d\ge0$.
Fix a smooth one-form $\beta$ with
$\Omega=\dd\beta$ symplectic and an $\Omega$-compatible smooth almost
complex structure $J$. Put $g=\Omega(\,\cdot\,,J\,\cdot\,)$.
Assume that $R:X\to[1,\infty)$ is smooth and proper, $c>0$, $r_*>1$,
and, on $\{R>r_*\}$,
\begin{equation}\label{eq:radial}
 \beta=-c\,\dd R\circ J,\qquad
 a=c|\dd R|_{g^*}^2\le R.
\end{equation}
The compact core $\{R\le r_*\}$ is exempt from these radial identities.

Let $\lambda$ be smooth, with $\omega=\dd\lambda$ symplectic and
\[
 \omega(v,Jv)\ge0\qquad(v\in TX).
\]
Assume that for some smooth $B:[1,\infty)\to(0,\infty)$ and $C>0$,
\begin{equation}\label{eq:growth}
 B(r)\ge r,\qquad
 \int_1^\infty\frac{r}{B(r)^2}\,\dd r=+\infty,\qquad
 \|\omega^{-1}\|_g\le C B(R)
\end{equation}
outside a compact set, with the first two conditions holding on
the whole half-line. The inverse norm is the operator norm of
$(v\mapsto\omega(v,\cdot))^{-1}$ from the $g$-dual to the $g$-tangent
space.

\begin{theorem}\label{thm:main}
There exist a diffeomorphism $F:X\to X$, smoothly isotopic to
the identity, and a globally smooth $h:X\to\R$ such that
\[
 F^*\lambda=\beta+\dd h,\qquad F^*\omega=\Omega.
\]
\end{theorem}
No bound on the primitive $\lambda$ is assumed. Neither completeness
of $g$ nor completeness of either Liouville field is used.

We construct positive radial functions $k$ and $A$, constant on
$[1,r_*]$, and compare the five primitives
\begin{equation}\label{eq:chain}
 \lambda\ \longrightarrow\ \lambda+k(R)\beta\
 \longrightarrow\ \lambda+A(R)\beta\
 \longrightarrow\ A(R)\beta\ \longrightarrow\ \beta .
\end{equation}
The actual forward paths used to construct the maps are
\begin{align}
 \mu^{(1)}_t&=\lambda+t\,k(R)\beta, &
 \mu^{(2)}_t&=\lambda+\bigl((1-t)k(R)+tA(R)\bigr)\beta,\nonumber\\
 \mu^{(3)}_t&=A(R)\beta+t\lambda, &
 \mu^{(4)}_t&=\bigl(1+t(A(R)-1)\bigr)\beta .
 \label{eq:fourpaths}
\end{align}
The third and fourth arrows in \eqref{eq:chain} use the inverses
of the corresponding forward endpoint maps. We retain this
orientation throughout the exactness calculation.

\section{Two pointwise geometric estimates}
\subsection{Weak taming and the resolvent contraction}
\begin{lemma}\label{lem:linear}
Let $\eta$ be a $J$-compatible symplectic form on a finite-dimensional
real vector space, and put $h=\eta(\,\cdot\,,J\,\cdot\,)$.
If a nondegenerate alternating form $\omega$ satisfies
$\omega(v,Jv)\ge0$, then $\omega+t\eta$ is nondegenerate for every
$t\ge0$. Moreover,
\[
 \iota_X(\omega+t\eta)=\iota_W\omega
 \quad\Longrightarrow\quad |X|_h\le |W|_h .
\]
\end{lemma}
\begin{proof}
Define $C$ by $\omega(v,w)=h(Cv,w)$ and put $T=-JC$.
Compatibility gives $\eta(v,w)=h(Jv,w)$, $J^*=-J$, and
\[
 h(Tv,v)=h(Cv,Jv)=\omega(v,Jv)\ge0.
\]
Consequently
\[
 |(T+t\id)v|_h^2
 =|Tv|_h^2+2t\,h(Tv,v)+t^2|v|_h^2
 \ge |Tv|_h^2+t^2|v|_h^2 .
\]
For $t>0$, this proves injectivity of $T+t\id$ and hence
invertibility by finite dimensionality. At $t=0$, invertibility
follows from nondegeneracy of $\omega$.
Since $C+tJ=J(T+t\id)$, this proves the first assertion.

The displayed inequality also gives
$\|T(T+t\id)^{-1}\|_h\le1$.
The operators $T$ and $(T+t\id)^{-1}$ commute, because $T$
commutes with $T+t\id$. The contraction equation is
$(T+t\id)X=TW$, and therefore
\[
 |X|_h=|(T+t\id)^{-1}TW|_h\le|W|_h.
\]
The argument uses no strictly positive lower bound for
$\omega(v,Jv)$.
\end{proof}

\subsection{The metric of a radial primitive}
Let $\rho:[1,\infty)\to(0,\infty)$ be smooth and constant on
$[1,r_*]$. Write
\[
 \eta_\rho=\dd(\rho(R)\beta),\qquad
 h_\rho(v,w)=\eta_\rho(v,Jw).
\]
On the core, $h_\rho=\rho(R)g$, because the coefficient derivative
vanishes there. Outside the core, the product rule and
\eqref{eq:radial} give
\begin{equation}\label{eq:metric}
 h_\rho(v,w)=\rho g(v,w)
 +c\rho'\bigl(\dd R(v)\dd R(w)+\dd R(Jv)\dd R(Jw)\bigr),
\end{equation}
where coefficients on $X$ are evaluated at $R$.

\begin{lemma}\label{lem:metric}
If $\rho(r)+r\rho'(r)>0$ for all $r\ge1$, then $\eta_\rho$
is $J$-compatible. On the exterior, put $D_\rho=\rho+\rho'a$.
Then
\begin{equation}\label{eq:dual}
 D_\rho>0,\qquad
 |\dd R|_{h_\rho^*}^2=\frac{a}{cD_\rho},\qquad
 |\beta|_{g^*}^2=ca.
\end{equation}
If $\rho'\le0$, then
$(\rho+R\rho')g\le h_\rho\le\rho g$.
If $\rho'\ge0$, then $h_\rho\ge\rho g$.
These metric inequalities hold on the whole manifold.
\end{lemma}
\begin{proof}
The right side of \eqref{eq:metric} is symmetric and invariant under
simultaneous application of $J$. At a point with $\dd R\ne0$, decompose
the tangent space into the $g$-orthogonal plane spanned by
$\nabla_g R,J\nabla_g R$ and its orthogonal complement.
The eigenvalues of $h_\rho$ relative to $g$ are $D_\rho$ on
that plane and $\rho$ on the complement.
If $\rho'\le0$, then $a\le R$ gives
$D_\rho\ge\rho+R\rho'>0$; if $\rho'\ge0$, then $D_\rho\ge\rho>0$.
At $\dd R=0$ the added tensor vanishes and $h_\rho=\rho g$.
These statements prove positivity and the metric inequalities.
The exterior derivative $\eta_\rho$ is closed, so positivity and
$J$-invariance prove compatibility and symplecticity.

The covector $\dd R$ is supported on the distinguished plane;
its squared dual norm is thus divided by $D_\rho$.
Since $|\dd R|_{g^*}^2=a/c$, this gives the first norm identity.
The second follows from $\beta=-c\,\dd R\circ J$ and the
$g$-orthogonality of $J$. Both identities hold at $\dd R=0$
as well. On the core, $\rho'=0$, which proves the remaining
global metric claims without using the exterior radial identities
there.
\end{proof}

\section{Constructing the two radial coefficients}
\subsection{A small coefficient adapted to the inverse-form bound}
Fix a smooth nondecreasing step function $\vartheta:\R\to[0,1]$
which is zero on $(-\infty,0]$ and one on $[1,\infty)$.
For example, with $q(t)=e^{-1/t}$ for $t>0$ and $q(t)=0$
otherwise, take $\vartheta(t)=q(t)/(q(t)+q(1-t))$.
Define
\[
 S(r)=1+\int_1^r\frac{u}{B(u)^2}\,\dd u,\qquad
 \chi(r)=\vartheta(r-r_*).
\]
Then $S\ge1$, $S'=r/B(r)^2>0$, $S(r)\to\infty$, and
$S(r)\le1+\log r$, using $B(r)\ge r$.
Set
\begin{equation}\label{eq:k}
 k(r)=\exp\left(\int_1^r
   \chi(s)\left(\frac{S'(s)}{2S(s)}-\frac1s\right)\dd s\right),
 \qquad m(r)=k(r)+rk'(r).
\end{equation}
Since $rS'(r)/(2S(r))\le1/2$, direct differentiation gives
\begin{align}
 &0<k\le1,\qquad k=1\text{ on }[1,r_*],\qquad k'\le0,\nonumber\\
 &m=k\left(1-\chi+\frac{\chi r^2}{2B(r)^2S(r)}\right)>0 .
 \label{eq:kproperties}
\end{align}
Strict positivity of the last expression holds both when $\chi<1$
and when $\chi=1$.

For $r\ge r_*+1$, the integrand in \eqref{eq:k} is
the derivative of $\frac12\log S-\log r$.
Thus there is $b_0>0$ such that
\begin{equation}\label{eq:ktail}
 k(r)=b_0\frac{\sqrt{S(r)}}r,\qquad
 m(r)=\frac{b_0r}{2B(r)^2\sqrt{S(r)}} .
\end{equation}
In particular $rk(r)\to\infty$.
Lemma~\ref{lem:metric} now supplies a compatible symplectic form
$\eta=\dd(k(R)\beta)$ and metric $h=\eta(\,\cdot\,,J\,\cdot\,)$
satisfying
\begin{equation}\label{eq:smallmetric}
 m(R)g\le h\le k(R)g.
\end{equation}
This construction works for the given arbitrary admissible $B$.

\subsection{A large coefficient with no upper growth restriction}
The continuity of the $g$-dual norm of $\lambda$ and properness of
$R$ give finite nondecreasing numbers
\[
 M_n=1+\sup_{\{R\le r_*+n+1\}}|\lambda|_{g^*}
 \qquad(n=0,1,2,\ldots).
\]
For an empty sublevel set, the displayed supremum is defined as
zero. Define
\begin{equation}\label{eq:A}
 A(r)=M_0+\sum_{j=1}^\infty
   (M_j-M_{j-1})\,\vartheta(r-r_*-j+1).
\end{equation}
On every bounded interval, all sufficiently large-index summands
vanish. The sum is therefore locally finite and smooth.
Every coefficient is nonnegative, so $A'\ge0$;
also $A=M_0$ on $[1,r_*]$.

If $r\in[r_*+n,r_*+n+1]$, each summand with $1\le j\le n$
has reached its full value, hence $A(r)\ge M_n$.
The sublevel defining $M_n$ contains every point with this value
of $R$. For $R\le r_*$, $M_0$ already controls the primitive.
Consequently
\begin{equation}\label{eq:Aproperties}
 A\ge1\ge k,\qquad A+rA'>0,\qquad
 A(R(x))\ge1+|\lambda_x|_{g^*}.
\end{equation}
There is no upper bound on the growth of $A$.
The interpolation argument below is designed to allow this freedom.

\section{Global time-dependent flows from proper barriers}
For a smooth path of primitives $\mu_t$ whose differentials are
symplectic, the Moser field is uniquely defined by
\begin{equation}\label{eq:moser}
 \iota_{X_t}\dd\mu_t=-\dot\mu_t .
\end{equation}
In local frames it is the inverse of a smooth invertible matrix
applied to a smooth covector. Hence it is jointly smooth on the
open nondegeneracy locus. When nondegeneracy holds on $[0,1]\times X$,
this locus supplies a neighborhood of the whole closed strip;
no uniform extra time interval on the noncompact manifold is
required.

\begin{lemma}[Compact confinement and continuation]\label{lem:continuation}
Let $X_t$ be such a smooth time-dependent field. If every solution
through a fixed initial time and point remains in a fixed compact
subset of $X$ as its time varies over its interval of existence
inside $[0,1]$, then that solution exists on the whole interval,
in both directions from the initial time.
If the confinement hypothesis holds for every initial datum,
the solution maps $E_{s,t}$ are smooth diffeomorphisms and satisfy
\[
 E_{s,s}=\id,\qquad
 E_{s,t}\circ E_{r,s}=E_{r,t},\qquad
 E_{s,t}^{-1}=E_{t,s}.
\]
They depend jointly smoothly on $s,t,x$ up to the endpoints.
\end{lemma}
\begin{proof}
Local existence, uniqueness, and smooth dependence follow in charts
from the smooth Euclidean ODE theorem. Around the compact spacetime
set $[0,1]\times K$ containing a confined curve, finitely many
local existence neighborhoods give a common positive lifetime
for nearby initial data.
If a maximal solution had a finite endpoint within the prescribed
time interval, restarting it sufficiently close to that endpoint
would extend it past the endpoint. Uniqueness identifies the new
solution with the old one on their overlap, a contradiction.
The argument is identical in reverse time.

Uniqueness gives the evolution identities and inverse maps.
Along any fixed finite trajectory, finitely many local flow maps
compose to give smooth dependence on both times and the initial
point. Equivalently, use the local flow of the spacetime field
$\partial_t+X_t$ on its open domain.
These local descriptions agree by uniqueness and give the asserted
smooth evolution, including one-sided endpoint derivatives and
local smooth extensions around the closed time square.
\end{proof}

\begin{corollary}[One proper scalar]\label{cor:scalar}
If a smooth proper function $P:X\to\R$ is bounded below and
$|X_tP|\le D$ on $[0,1]\times X$, then the conclusions of
Lemma~\ref{lem:continuation} hold.
An estimate valid outside one fixed compact set is sufficient.
\end{corollary}
\begin{proof}
Along a trajectory starting at $(s,x)$,
$|P(\gamma(t))-P(x)|\le D|t-s|$.
Its image lies in the compact inverse image of a bounded closed
interval. If the estimate is initially exterior, continuity of
$(t,x)\mapsto X_tP(x)$ on the product of the compact time interval
and the exceptional compact set gives one enlarged global bound.
\end{proof}
This uses properness, not metric completeness. In particular, a
bound on the norm of a field is used below only after conversion
to a bound on the derivative of a proper scalar.

\section{The interpolation lemma and its reverse-time barrier}
Let $\sigma$ be a smooth one-form with
$\dd\sigma(v,Jv)\ge0$; its differential may be degenerate.
Let $\rho_0,\rho_1$ be positive smooth functions, constant on
$[1,r_*]$, with
\begin{equation}\label{eq:bridgehyp}
 \rho_i+r\rho_i'>0,\qquad
 \rho_1\ge\rho_0,\qquad r\rho_0(r)\longrightarrow\infty.
\end{equation}
Put $\delta=\rho_1-\rho_0\ge0$ and
$\rho_t=(1-t)\rho_0+t\rho_1$.

\begin{lemma}[Radial interpolation]\label{lem:bridge}
The path $\mu_t=\sigma+\rho_t(R)\beta$ has symplectic differential.
Its Moser evolution exists for all pairs of times in $[0,1]$,
with the smoothness, inverse, and evolution properties of
Lemma~\ref{lem:continuation}.
\end{lemma}
\begin{proof}
The coefficient inequalities persist under convex combinations.
Define the compatible metric
\[
 h_t=\dd(\rho_t(R)\beta)(\,\cdot\,,J\,\cdot\,).
\]
Lemma~\ref{lem:metric} and weak taming give
\[
 \dd\mu_t(v,Jv)=\dd\sigma(v,Jv)+|v|_{h_t}^2>0\quad(v\ne0),
\]
which proves nondegeneracy. Closedness is automatic.
The Moser equation is $\iota_{X_t}\dd\mu_t=-\delta(R)\beta$.

On the exterior, evaluate it on $JX_t$:
\begin{equation}\label{eq:bridgeenergy}
 |X_t|_{h_t}^2+\dd\sigma(X_t,JX_t)
   =-c\delta\,X_tR .
\end{equation}
If $\delta=0$, nondegeneracy implies $X_t=0$ at the point.
Otherwise \eqref{eq:bridgeenergy} gives $Y:=-X_tR\ge0$.
Using \eqref{eq:dual}, with $D_t=\rho_t+\rho_t'a$, yields
\[
 Y^2\le|\dd R|_{h_t^*}^2|X_t|_{h_t}^2
 \le\frac{a}{cD_t}\,c\delta Y .
\]
For $Y>0$ divide by $Y$; for $Y=0$ the result is immediate.
Thus
\begin{equation}\label{eq:radialdrift}
 0\le-X_tR\le\frac{\delta a}{D_t}
 \quad\text{on }\{R>r_*\}.
\end{equation}

Define the moving scalar $Q(t,x)=R(x)\rho_t(R(x))$.
Its total derivative along the field satisfies, on the exterior,
\begin{align}
 (\partial_t+X_t)Q
 &=\delta R+(\rho_t+R\rho_t')X_tR\nonumber\\
 &\ge\delta\left(R-\frac{a(\rho_t+R\rho_t')}{D_t}\right)
 =\frac{\delta\rho_t(R-a)}{D_t}\ge0 .
 \label{eq:movingbarrier}
\end{align}
The sign of $\rho_t+R\rho_t'$ is positive, so using the lower
bound for $X_tR$ in this calculation is legitimate.

Compactness of the core and the time interval supplies constants
$C_0,C_1\ge0$ such that, globally,
\[
 X_tR\le C_0,\qquad (\partial_t+X_t)Q\ge-C_1 .
\]
For a trajectory with $\gamma(s)=x$ and $t\ge s$,
$R(\gamma(t))\le R(x)+C_0(t-s)$.
For $t\le s$, integrate the other inequality to obtain
\[
 Q(t,\gamma(t))\le Q(s,x)+C_1(s-t).
\]
But $Q(t,y)\ge R(y)\rho_0(R(y))$, and the latter expression
tends to infinity with $R(y)$ by \eqref{eq:bridgehyp}.
This also bounds $R(\gamma(t))$, now in reverse time.
Both portions of the trajectory lie in a single compact
$R$-sublevel, so Lemma~\ref{lem:continuation} applies.
\end{proof}

\begin{remark}[Why the two barriers matter]
The argument does not assert a bound on $|X_tR|$ independent of
$R$. The forward bound controls $R$, while the reverse bound controls
$Q$. Arbitrarily fast growth of the upper coefficient $\rho_1$
does not spoil the lower estimate
$Q(t,x)\ge R(x)\rho_0(R(x))$.
If one constructs trajectories using compactly supported cutoffs,
these estimates remain available: for $0\le\zeta\le1$,
\[
 (\partial_t+\zeta X_t)Q
 =\zeta(\partial_t+X_t)Q+(1-\zeta)\partial_tQ,
\]
and $\partial_tQ=\delta(R)R\ge0$. Thus the cutoff does not
destroy the reverse-time inequality.
\end{remark}

\section{The first path: the integral growth hypothesis}
Consider $\mu_t=\lambda+t\,k(R)\beta$.
Lemma~\ref{lem:linear}, applied to $\eta=\dd(k(R)\beta)$,
proves that $\dd\mu_t=\omega+t\eta$ is symplectic.
Define $W$ by $\iota_W\omega=-k(R)\beta$.
The actual Moser field satisfies
\[
 \iota_{X_t}(\omega+t\eta)=\iota_W\omega,\qquad |X_t|_h\le|W|_h .
\]
On the region where the inverse-form bound holds,
\[
 |W|_g\le C B(R)k(R)|\beta|_{g^*}.
\]
The metric comparison \eqref{eq:smallmetric} then gives
\[
 |X_t|_g\le\sqrt{\frac{k(R)}{m(R)}}\,|W|_g .
\]
On the radial exterior,
$|\dd R|_{g^*}|\beta|_{g^*}=a\le R$ by \eqref{eq:dual}.
For $R\ge r_*+1$, outside the original exceptional compact set,
\begin{align}
 |X_tR|
 &\le C B(R)k(R)R\sqrt{\frac{k(R)}{m(R)}}\nonumber\\
 &=\sqrt2\,Cb_0\,\frac{B(R)^2S(R)}R .
 \label{eq:firstdrift}
\end{align}
Here the tail formulas \eqref{eq:ktail} were used exactly.

Now set $P=\log S(R)$. Since $S\ge1$ and $S(r)\to\infty$,
this is a smooth proper function bounded below by zero.
For example, a bounded upper value of $\log S(R)$ bounds $R$,
and closed sublevels are compact by properness of $R$.
Using $S'=r/B(r)^2$ in \eqref{eq:firstdrift} gives
\[
 |X_tP|\le\sqrt2\,Cb_0 .
\]
The region excluded from this estimate is contained in the union
of the original compact exception and the compact set
$\{R\le r_*+1\}$. Smoothness supplies a bound there too.
Corollary~\ref{cor:scalar} therefore constructs the full
two-time smooth evolution of the first path.

This is the only stage whose completeness estimate uses the
inverse-form growth bound. Its right-hand side is $-k(R)\beta$,
so no size bound on $\lambda$ enters the argument.

\section{The other three paths}
\subsection{The second path}
Apply Lemma~\ref{lem:bridge} with
\[
 \sigma=\lambda,\qquad \rho_0=k,\qquad \rho_1=A.
\]
Weak taming gives the hypothesis on $\dd\sigma$.
Equations \eqref{eq:kproperties}, \eqref{eq:ktail}, and
\eqref{eq:Aproperties} give all coefficient hypotheses, including
$rk(r)\to\infty$. This constructs the smooth evolution for
$\mu^{(2)}$ in \eqref{eq:fourpaths}. Its reverse-time barrier is
bounded below by $Rk(R)$ independently of how quickly $A$ grows.

\subsection{The third path}
Use the forward path $\mu_s=A(R)\beta+s\lambda$.
Since $A'\ge0$, Lemma~\ref{lem:metric} gives
\[
 \dd\mu_s(v,Jv)
 =\dd(A(R)\beta)(v,Jv)+s\omega(v,Jv)
 \ge A(R)|v|_g^2 .
\]
Thus every differential is symplectic, including at $s=0$.
The Moser equation is $\iota_{X_s}\dd\mu_s=-\lambda$.
At $JX_s$ it implies
\[
 A(R)|X_s|_g^2
 \le-\lambda(JX_s)
 \le|\lambda|_{g^*}|X_s|_g .
\]
If $X_s=0$, the following bound is immediate; otherwise divide by
its positive norm. Equation \eqref{eq:Aproperties} yields
\begin{equation}\label{eq:thirdbound}
 |X_s|_g\le\frac{|\lambda|_{g^*}}{A(R)}\le1 .
\end{equation}

Use the proper smooth scalar $P=\sqrt R$, rather than a
metric completeness assertion. On the radial exterior,
\[
 |X_sP|
 \le\frac{|\dd R|_{g^*}}{2\sqrt R}|X_s|_g
 \le\frac1{2\sqrt c}.
\]
The derivative is bounded on the core by joint continuity and
compactness. Corollary~\ref{cor:scalar} constructs the full
evolution in both directions.
The forward endpoint goes from the primitive $A(R)\beta$
to $A(R)\beta+\lambda$ in the Moser pullback sense.
Its inverse supplies the third arrow of \eqref{eq:chain}.

\subsection{The fourth path}
Apply Lemma~\ref{lem:bridge} with
$\sigma=0$, $\rho_0=1$, and $\rho_1=A$.
The lower divergence condition is simply $r\to\infty$.
It constructs the evolution for
$\mu^{(4)}_t=(1+t(A(R)-1))\beta$, from $\beta$ to $A(R)\beta$.
The inverse endpoint supplies the fourth arrow.

For each of the four paths let $G_{j,t}=E^{(j)}_{0,t}$,
and write $G_j=G_{j,1}$. Every $G_{j,t}$ and its inverse depend
smoothly on $(t,x)$, with $G_{j,0}=\id$. All four endpoint maps
are therefore smoothly isotopic to the identity, as are their
inverses.

\section{The exact pullback identity and its smooth gauge}
We prove the identity for any of the smooth Moser paths above.
Write $\mu_t$, $X_t$, and $G_t$ for its primitive, field, and
identity-start evolution.

\begin{lemma}[Time-dependent pullback formula]\label{lem:cartan}
For $0\le t\le1$,
\[
 \frac{\dd}{\dd t}G_t^*\mu_t
 =G_t^*(\dot\mu_t+\iota_{X_t}\dd\mu_t)
   +\dd\bigl(\mu_t(X_t)\circ G_t\bigr).
\]
Derivatives at the two endpoints are interpreted within $[0,1]$.
\end{lemma}
\begin{proof}
Fix $x$ and $v\in T_xX$. In local coordinates along the
trajectory, put $y=G_t(x)$ and $w=\dd G_t(x)v$.
Differentiating $\mu_t(y)(w)$ gives three terms:
\[
 \dot\mu_t(y)(w)
 +(D\mu_t(y)[X_t(y)])(w)
 +\mu_t(y)\bigl(DX_t(y)[w]\bigr).
\]
The equation for the differentiated trajectory,
$\partial_t w=DX_t(y)[w]$, follows by differentiating
$\partial_tG_t=X_t\circ G_t$ in the initial point and commuting
mixed derivatives of the smooth local expression for $G$.
For a one-form in coordinates,
\[
 \dd\mu_t(X_t,w)
 =(D\mu_t[X_t])(w)-(D\mu_t[w])(X_t).
\]
Meanwhile the derivative in direction $v$ of
$\mu_t(X_t)\circ G_t$ is
$(D\mu_t[w])(X_t)+\mu_t(DX_t[w])$.
Combining these formulas proves the identity.
Each expression is intrinsic, so the local computations agree
on chart overlaps. Smoothness up to the endpoints supplies the
corresponding one-sided identities there.
\end{proof}

Substituting the actual Moser equation \eqref{eq:moser} cancels
the first two terms:
\begin{equation}\label{eq:derivativeexact}
 \frac{\dd}{\dd t}G_t^*\mu_t
 =\dd\bigl(\mu_t(X_t)\circ G_t\bigr).
\end{equation}

\begin{lemma}[Smooth parameter integration]\label{lem:integral}
If $f:[0,1]\times X\to\R$ is jointly smooth up to the endpoints,
then $H(x)=\int_0^1 f(t,x)\,\dd t$ is smooth, and
\[
 \dd H_x(v)=\int_0^1\dd_x f(t,x)(v)\,\dd t .
\]
\end{lemma}
\begin{proof}
Work in a coordinate chart about a fixed point $x_0$.
Smoothness up to the endpoints gives local smooth extensions.
Compactness of $[0,1]\times\{x_0\}$ permits finitely many such
neighborhoods and then a smaller spatial neighborhood on which
the relevant derivatives are defined throughout the time interval.
Choose its closure inside a relatively compact coordinate
neighborhood. Each fixed finite-order spatial derivative is
continuous and hence bounded on the resulting compact product.

For the first derivative, apply the fundamental theorem of calculus
to the spatial difference quotient. Continuity of the derivative
on that compact product gives uniform convergence as the increment
tends to zero, allowing integration in $t$ to pass to the limit.
This gives the displayed derivative formula.
Repeating the same argument for every finite order proves
$C^\infty$ regularity. Coordinate changes obey the usual chain
rule, so the local smooth functions and their differentials
give the stated global result.
\end{proof}
Only local bounds near each spatial point are needed.
There is no global bound on derivatives of $f$ over the noncompact
manifold and no assumption of a uniform extra time strip.
Assigning arbitrary values to $f$ outside $[0,1]$ when writing
the integral does not assert smoothness of that arbitrary extension.

\begin{proposition}[Exactness of one endpoint]\label{prop:exactness}
Define the actual scalar function
\[
 H(x)=\int_0^1\mu_t(X_t)(G_t(x))\,\dd t .
\]
Then $H$ is smooth and
\begin{equation}\label{eq:exactendpoint}
 G_1^*\mu_1=\mu_0+\dd H,\qquad
 G_1^*\dd\mu_1=\dd\mu_0.
\end{equation}
\end{proposition}
\begin{proof}
The integrand is jointly smooth by smoothness of the form,
the vector field, and the evolution. Lemma~\ref{lem:integral}
proves smoothness of $H$ and the derivative-interchange formula.
For fixed $x,v$, integrate \eqref{eq:derivativeexact} from zero
to one and use $G_0=\id$. The fundamental theorem of calculus
then gives the primitive equality evaluated on $v$.
Taking exterior derivatives, using naturality of $\dd$ under
pullback and $\dd^2H=0$, gives the two-form equality.
\end{proof}
Thus the gauge is constructed by the displayed integral; its
existence or smoothness is not an additional premise.

\section{Composition, inversion, and the final isotopy}
For clarity, write the five primitives in \eqref{eq:chain} as
\[
 \nu_0=\lambda,\quad \nu_1=\lambda+k(R)\beta,\quad
 \nu_2=\lambda+A(R)\beta,\quad
 \nu_3=A(R)\beta,\quad \nu_4=\beta .
\]
Applying Proposition~\ref{prop:exactness} to the four forward paths
gives smooth $H_1,H_2,H_3,H_4$ with
\begin{equation}\label{eq:fouridentities}
\begin{aligned}
 G_1^*\nu_1&=\nu_0+\dd H_1,&
 G_2^*\nu_2&=\nu_1+\dd H_2,\\
 G_3^*\nu_2&=\nu_3+\dd H_3,&
 G_4^*\nu_3&=\nu_4+\dd H_4 .
\end{aligned}
\end{equation}

\begin{lemma}[Algebra of exact comparisons]\label{lem:algebra}
If $F^*\alpha=\beta+\dd h$ and $G^*\beta=\gamma+\dd k$, then
\[
 (F\circ G)^*\alpha=\gamma+\dd(k+h\circ G).
\]
If $F^*\alpha=\beta+\dd h$, then
\[
 (F^{-1})^*\beta=\alpha+\dd(-h\circ F^{-1}).
\]
\end{lemma}
\begin{proof}
The first identity follows by pulling the first equality back by
$G$, substituting the second equality, and using
$G^*\dd h=\dd(h\circ G)$.
For the second, pull the supplied identity back by $F^{-1}$
and rearrange. These calculations also fix the signs of the
gauges under inversion.
\end{proof}

Set
\[
 P_2=G_2\circ G_1,\qquad
 P_3=G_3^{-1}\circ P_2,\qquad
 \Psi=G_4^{-1}\circ P_3.
\]
Successive applications of Lemma~\ref{lem:algebra} to
\eqref{eq:fouridentities} yield
\[
 \Psi^*\beta=\lambda+\dd f,\qquad
 f=H_1+H_2\circ G_1-H_3\circ P_3-H_4\circ\Psi .
\]
Therefore the witnesses for Theorem~\ref{thm:main} are
\begin{equation}\label{eq:final}
 F=\Psi^{-1}=G_1^{-1}\circ G_2^{-1}\circ G_3\circ G_4,
 \qquad h=-f\circ F.
\end{equation}
They satisfy $F^*\lambda=\beta+\dd h$, and exterior differentiation
gives $F^*\omega=\Omega$.

An explicit identity isotopy is available without choosing
junctions between separate time intervals:
\[
 F_t=G_{1,t}^{-1}\circ G_{2,t}^{-1}\circ G_{3,t}\circ G_{4,t}.
\]
Every factor and its inverse is jointly smooth.
The inverse family is the reverse-order composition of the inverse
factors and is jointly smooth as well.
Since every $G_{j,0}=\id$, this family starts at the identity and
ends at the $F$ in \eqref{eq:final}.
This proves the required smooth isotopy and completes the theorem.

\section{Cylindrical compatible structures}
We prove the cylindrical assertion in the companion statement. Fix its
Liouville domain, completion, boundary $Y$, and form $\alpha$.
The boundary restriction of the symplectic form is $\dd\alpha$.
Outward transversality of the Liouville field makes $\alpha$ contact:
the restriction of $\dd\alpha$ to $\xi=\ker\alpha$ is nondegenerate.
The normalized Reeb field is the unique solution of
$\alpha(R_\alpha)=1$ and $\iota_{R_\alpha}\dd\alpha=0$.
In local frames this is an invertible linear system with smooth
coefficients, so its solutions form a smooth field. The constant-rank
kernel $\xi$ is an actual smooth vector subbundle of $TY$.

\subsection{Smooth symplectic normalization}
Let $\Omega$ be a smooth nondegenerate alternating form on a real
vector bundle and let $h$ be a smooth positive metric. Define $A$ by
\[
 h(Au,v)=\Omega(u,v).
\]
Then $A$ is invertible and $h$-skew-adjoint. The operator
$S=-A^2=A^*A$ is positive and self-adjoint. Put
\[
 P=S^{1/2},\qquad J=AP^{-1}.
\]
The positive square root commutes with $A$, hence $J^2=-\id$.
Furthermore,
\[
 \Omega(u,Jv)=h(Au,AP^{-1}v)=h(Pu,v),
\]
which is symmetric and strictly positive for $u=v\ne0$.
Since $J$ is $h$-orthogonal and commutes with $A$, it preserves
$\Omega$ as well.

These constructions are smooth in the base point. To see the only
nonpolynomial step, in local smooth orthonormal frames the derivative
of the square map at a positive matrix $P$ is
$H\mapsto PH+HP$. In an eigenbasis it multiplies the $(i,j)$ entry
by $p_i+p_j>0$, so it is invertible. The implicit function theorem
gives local smooth positive square roots; uniqueness makes them agree
on overlaps. Inversion is smooth on invertible matrices. This also
gives the intrinsic construction independently of the chosen frames.

There is an exact recovery property. If $J_0$ is already compatible
and $h(u,v)=\Omega(u,J_0v)$, then
$h(J_0u,v)=\Omega(u,v)$, so $A=J_0$, $P=\id$, and the construction
returns $J_0$. This equality, not just a homotopy class, will preserve
the prescribed end.

Applied to $\xi$ with $\Omega=\dd\alpha|_\xi$, the construction
produces at least one smooth compatible $j_\xi$: choose a smooth
positive metric by gluing local chart metrics with a partition of
unity, then normalize. Smoothness of the native restriction and
inclusion identifies this bundle construction with the literal
contact distribution. Empty bases and zero-rank fibers cause no
exception to the construction.

\subsection{The prescribed operator on a neighborhood of the end}
Now fix any prescribed smooth compatible $j_\xi$. Split a product
tangent vector as
\[
 (v,s)=bR_\alpha+\eta+s\partial_t,\qquad
 b=\alpha(v),\quad\eta\in\xi,
\]
and set
\[
 J_{\rm end}(v,s)=(j_\xi\eta+sR_\alpha,-b).
\]
Its square is minus the identity. On product coordinates the actual
ambient form pulls back to
\[
 E^*\dd\beta=e^t(\dd t\wedge\alpha+\dd\alpha).
\]
The Reeb annihilation identity and compatibility of $j_\xi$ give
\[
 E^*\dd\beta\bigl((v,s),J_{\rm end}(v,s)\bigr)
 =e^t\bigl(s^2+b^2+\dd\alpha(\eta,j_\xi\eta)\bigr)>0
\]
for nonzero $(v,s)$. The same decomposition proves preservation of
the two-form and symmetry of the induced pairing.

Use the completion's flow and end immersion to transport this operator
to its open escaping basin. Uniqueness of the flow and outward
transversality give the native basin coordinates; the differential
and its inverse identify the actual tangent bundles. Thus the
transported operator and
$h_{\rm end}=\dd\beta(\,\cdot\,,J_{\rm end}\,\cdot\,)$ are smooth
on an open neighborhood of the entire closed end
$K=E(Y\times[0,\infty))$. In particular, this neighborhood extends
past the seam; a formula defined only for $t>0$ would not suffice.
The completion geometry proves that $K$ is closed and lies in that
open basin.

\subsection{Relative extension and recovery on the whole closed end}
We use the following relative metric fact: if $K$ is closed, $U$ is
open with $K\subset U$, and a smooth positive metric $q$ is given
on $U$, then a global smooth positive metric $h$ exists with $h=q$
on every point of $K$. Here $K$ need not be compact.

For completeness, use the convex set of positive symmetric forms in
each tangent fiber, additionally requiring equality to $q$ over $K$.
Near a point of $K$, the prescribed $q$ supplies a local section.
Near a point outside $K$, choose a chart neighborhood disjoint from
$K$ and use a positive chart metric. A locally finite smooth
partition of unity glues these sections. Positivity is preserved by
convex combinations, and every contributing section equals $q$ at
points of $K$. This proves the required equality also at the boundary
of $K$. Finite-dimensional native fiber topology gives the bounded
unit-ball property required to register the result as a Riemannian
metric.

Apply this fact to $h_{\rm end}$ and normalize the resulting global
$h$ against $\Omega=\dd\beta$. Smooth normalization gives a global
compatible $J$, and the recovery property gives exactly $J_{\rm end}$
on all of $K$. Define the output metric by
$g=\dd\beta(\,\cdot\,,J\,\cdot\,)$; it is smooth and positive by
the preceding calculation. The auxiliary $h$ is not identified with
$g$ away from $K$.

This proves arbitrary prescribed extension, including $t=0$.
Supplying the contact structure constructed above proves unconditional
cylindrical existence. The same construction works on the quotient
supplied by the native finite-type adapter: use its domain, completed
primitive, and end map throughout. Its source identifications remain
unchanged. The comparison application then chooses this $J$ before
quantifying the weakly tamed target form and retains the original
growth hypothesis.

\section{The finite-type Liouville application}
Use the supplied domain and smooth completion realization from
Section~4 of the companion \emph{Definitions, hypotheses, and target
theorems}: the compact-piece inclusion $j$, complete reference flow
$\varphi_t$, and actual end immersion $E$ are part of the data.
In particular, $E(y,t)=\varphi_t(j(b(y)))$, $E^*\beta=e^t\alpha$,
and the embedded nonnegative end together with $j(W)$ covers $X$,
with positive times disjoint from $j(W)$. Write $r=e^t$ on this end.
Then $\beta=r\alpha$ and $\partial_t=r\partial_r$.
The preceding construction lets us choose a compatible cylindrical $J$ with
$J(r\partial_r)=R_\alpha$ and an $r$-independent compatible
structure on $\ker\alpha$.

On the end, decompose a tangent vector as
$v=a\,r\partial_r+bR_\alpha+\xi$ with $\xi\in\ker\alpha$.
Then
$Jv=aR_\alpha-b\,r\partial_r+J\xi$, and hence
$-\dd r(Jv)=br=\beta(v)$.
Moreover,
\[
 \dd\beta=\dd r\wedge\alpha+r\,\dd\alpha,\qquad
 g(r\partial_r,r\partial_r)=r .
\]
The radial line is orthogonal to the Reeb and contact directions.
It follows that $\nabla_g r=r\partial_r$, so
\[
 \beta=-\dd r\circ J,\qquad |\dd r|_{g^*}^2=r .
\]

Choose a smooth cutoff $\chi$ on $[1,\infty)$ equal to zero
near $1$ and one for $r\ge2$, with values in $[0,1]$.
Define $R=1+\chi(r)(r-1)$ on the end and $R=1$ on the compact
piece. It glues smoothly, is at least one, equals $r$ for $r\ge2$,
and is at most two on the remaining compact region.
It is proper because $Y$ and $W$ are compact and $R=r$ far out.
Thus the radial hypotheses hold with $c=1$ and $r_*=2$.

If $\omega=\dd\lambda$ weakly tames this $J$ and obeys the
inverse-form estimate $CB(r)$ for sufficiently large $r$,
it obeys the main theorem's estimate $CB(R)$ outside a compact
set. Theorem~\ref{thm:main} gives the required exact comparison.
The reference completion already carries its prescribed complete
all-real-time Liouville flow; on the end its field is
$r\partial_r$, whose logarithmic-radius coordinate has unit speed.
No completeness of the other primitive
is used in this application.

\section{The properly embedded Stein application}
Let $X$ be a smooth real manifold, let $J$ be a smooth almost complex
structure, and let $\iota:X\hookrightarrow\C^N$ be a smooth proper
topological embedding with injective differential and
$\dd\iota\circ J=I\circ\dd\iota$.
These are the explicit input data of the companion statement; no
separate complex atlas or integrability premise is used in the proof.
Write the ambient coordinates as $z_\nu=x_\nu+iy_\nu$ and define
\[
 \beta_{\mathrm{E}}=\frac12\sum_{\nu=1}^N
   (x_\nu\,\dd y_\nu-y_\nu\,\dd x_\nu),\qquad
 \Omega_{\mathrm{E}}=\dd\beta_{\mathrm{E}}
   =\sum_{\nu=1}^N\dd x_\nu\wedge\dd y_\nu .
\]
If $I$ is the standard complex structure, then
$\beta_{\mathrm{E}}=-\frac14\dd\|z\|^2\circ I$ and
$\Omega_{\mathrm{E}}(u,Iv)=\ip{u}{v}_{\R}$.

Put $R=e+\|\iota\|^2$.
The assumed differential identity gives
\begin{equation}\label{eq:steinpullback}
 \beta:=-\frac14\dd R\circ J=\iota^*\beta_{\mathrm{E}},
 \qquad \Omega=\dd\beta=\iota^*\Omega_{\mathrm{E}}.
\end{equation}
Consequently
\[
 \Omega(v,Jw)=
 \Omega_{\mathrm{E}}(\dd\iota(v),I\,\dd\iota(w))
 =\ip{\dd\iota(v)}{\dd\iota(w)}_{\R}.
\]
Because $\iota$ is an immersion, this is positive definite.
It is smooth in $x$, so it is exactly the smooth induced
Euclidean metric $g$. The identity also proves nondegeneracy
of $\Omega$; closedness follows from \eqref{eq:steinpullback}.
Compatibility with $J$ follows from the corresponding ambient
identities. Thus the reference is determined from the embedding.

Properness of $\iota$ implies properness of $R$: a bounded
$R$-sublevel is the inverse image of a closed Euclidean ball,
and hence compact.
For every tangent vector,
\[
 |\dd R(v)|=2|\ip{\iota(x)}{\dd\iota(v)}_{\R}|
 \le2\|\iota(x)\|\,|v|_g .
\]
Taking the $g$-dual norm gives
\[
 |\dd R|_{g^*}^2\le4\|\iota(x)\|^2=4(R-e).
\]
Therefore the radial identity holds everywhere with $c=1/4$,
and $c|\dd R|_{g^*}^2\le R-e<R$.
Any fixed $r_*>1$ may be used for the main theorem's exterior
formulation.

It remains to align the lower endpoint of the growth function.
Suppose $B$ is admissible on $[e,\infty)$.
Choose a smooth step $\zeta$ equal to zero for $r\le e+1$
and one for $r\ge e+2$.
Define
\[
 \widetilde B(r)=
 \begin{cases}
 r,&1\le r\le e+1,\\
 r+\zeta(r)(B(r)-r),&r>e+1.
 \end{cases}
\]
Smoothness of the step at its constant region makes this smooth
across the join. It is positive and at least $r$, and equals
$B$ for $r\ge e+2$.
Changing the integrand on a bounded interval preserves divergence
of $\int r/\widetilde B(r)^2\,\dd r$.
Enlarging the exceptional compact set by $\{R\le e+2\}$
preserves the inverse-form bound.
Theorem~\ref{thm:main} now applies with $\widetilde B$.

This proves the Stein-reference corollary without a finite-type
hypothesis. The geometric verification concerns the given
embedding and its induced reference; it does not require choosing
a new Stein manifold or proving an embedding existence theorem.

\section{The common-reference consequence}
Suppose $\lambda_1$ and $\lambda_2$ both satisfy the main theorem
for one fixed radial reference. Apply the theorem separately,
retaining their independent growth witnesses, to obtain
\[
 F_i^*\lambda_i=\beta+\dd h_i\qquad(i=1,2).
\]
Pull the first equality back by $F_2^{-1}$ and use the inverse
of the second equality. The result is
\[
 (F_1\circ F_2^{-1})^*\lambda_1
 =\lambda_2+
   \dd\bigl((h_1-h_2)\circ F_2^{-1}\bigr).
\]
Taking exterior derivatives gives the symplectic equality.
Composing the identity isotopy of $F_1$ with the inverse isotopy
of $F_2$ gives the required smooth isotopy of the same comparison
map. Thus both primitives are compared with the same reference,
while their constants, functions $B_i$, and compact exceptions
remain independent.

\section{Where the hypotheses enter}
\begin{tabularx}{\textwidth}{@{}p{0.30\textwidth}X@{}}
\toprule
Hypothesis & Role in the proof\\
\midrule
Compatibility and explicit $g$ &
Give the radial metric formula, dual norms, and the orthogonality
of $J$ used in all estimates.\\
Weak, nonstrict taming &
Gives the resolvent contraction for the first path and positivity
after adding a compatible radial form for the other paths.\\
Properness of $R$ &
Bounds the primitive on each compact sublevel and converts
scalar trajectory bounds into compact confinement.\\
Radial identities only on the exterior &
Give the exterior estimates; compactness supplies the missing
bounds on the unrestricted core.\\
Inverse-form bound and divergent integral &
Construct $k$ and the proper first-path scalar $\log S(R)$.\\
No primitive growth bound &
The coefficient $A$ is constructed to dominate the actual
primitive, with no upper growth restriction.\\
Both time directions &
Produce actual inverse endpoint maps; reverse interpolation
uses the moving barrier $Q$ rather than a forward estimate alone.\\
Joint smoothness &
Gives smooth diffeomorphisms, the Cartan calculation, the
integral gauge, and the final identity isotopy.\\
\bottomrule
\end{tabularx}

\medskip
All four flows are finite-parameter Moser evolutions.
The general theorem does not turn their existence into a
claim that the autonomous Liouville fields of the prescribed
primitives are complete. Its conclusion is precisely the exact
primitive identity, its exterior-derivative consequence, and
the smooth identity isotopy stated at the beginning.
\end{document}
